% ============================================================
\chapter{Limit Cycles and the Poincar\'e-Bendixson Theorem}
\label{ch:limit_cycles}
% ============================================================

The heart beats. An electronic circuit oscillates. A predator-prey population fluctuates periodically. These phenomena of self-sustained oscillation are ubiquitous in nature and technology. Mathematically, they correspond to \emph{limit cycles}: isolated periodic orbits toward which neighbouring trajectories converge (or diverge). The Poincar\'e-Bendixson theorem (1901) is the fundamental result guaranteeing the existence of such cycles in planar systems: if a trajectory remains confined to a compact region with no equilibrium point, it must converge to a limit cycle. This result, specific to dimension 2, has no analogue in higher dimensions---and that is where chaos enters the picture.

\section{Introduction}

Limit cycles are isolated periodic orbits in phase space --- closed orbits
that do not belong to a continuous family of closed orbits. They constitute
one of the most important phenomena unique to nonlinear systems, since linear
systems can only produce centers (continuous families of closed orbits). The
Poincar\'e-Bendixson theorem provides a powerful criterion for proving the
existence of limit cycles in the plane.

\section{Definitions}

\begin{definition}[Limit cycle]
A \textbf{limit cycle} is a periodic orbit $\Gamma$ that is the
$\omega$-limit set or $\alpha$-limit set of at least one orbit not
belonging to $\Gamma$.
\end{definition}

\begin{definition}[Types of limit cycles]
A limit cycle $\Gamma$ is:
\begin{itemize}
    \item \textbf{stable} (attracting) if $\Gamma = \omega(x_0)$ for all
    $x_0$ in a neighborhood of $\Gamma$;
    \item \textbf{unstable} (repelling) if $\Gamma = \alpha(x_0)$ for all
    $x_0$ in a neighborhood of $\Gamma$;
    \item \textbf{semi-stable} if $\Gamma$ is attracting on one side and
    repelling on the other.
\end{itemize}
\end{definition}

\begin{example}[Limit cycle in the Van der Pol oscillator]
The Van der Pol equation
\[
    \ddot{x} - \mu(1 - x^2)\dot{x} + x = 0, \qquad \mu > 0,
\]
possesses a unique stable limit cycle. For small $\mu$, this cycle is close
to the circle of radius 2 centered at the origin. For large $\mu$, the cycle
takes on a characteristic \textbf{relaxation oscillation} shape.
\end{example}

\section{Non-existence Criterion: Bendixson-Dulac}

\begin{theorem}[Bendixson criterion]
\label{thm:bendixson}
Let $D \subseteq \R^2$ be a simply connected domain. If
$\mathrm{div}\, f(x) = \frac{\partial f_1}{\partial x_1} + \frac{\partial f_2}{\partial x_2}$
has constant sign (and is not identically zero) on $D$, then the system
$\dot{x} = f(x)$ has \textbf{no periodic orbit} entirely contained in $D$.
\end{theorem}

\begin{proof}
By Green's theorem, if $\Gamma$ were a periodic orbit enclosing a domain
$\Omega$, we would have
\[
    \iint_\Omega \mathrm{div}\, f\, dx_1\, dx_2 =
    \oint_\Gamma (f_1\, dx_2 - f_2\, dx_1) = 0,
\]
since on $\Gamma$, $dx_i = f_i\, dt$, and the line integral simplifies.
But if $\mathrm{div}\, f$ has constant sign and is not identically zero, the
double integral is nonzero: contradiction.
\end{proof}

\begin{theorem}[Dulac criterion]
The same result holds if $\mathrm{div}\, f$ is replaced by
$\mathrm{div}(\rho f)$ where $\rho : D \to \R$ is a strictly positive
$C^1$ function (Dulac function).
\end{theorem}

\begin{example}[Lotka-Volterra system]
For $\dot{x} = x(\alpha - \beta y)$, $\dot{y} = y(\delta x - \gamma)$ in the
first quadrant, taking $\rho(x,y) = \frac{1}{xy}$ gives
$\mathrm{div}(\rho f) = 0$, which is inconclusive. However, the existence
of a first integral shows that orbits are closed (center) and not limit cycles.
\end{example}

\section{The Poincar\'e-Bendixson Theorem}

\begin{theorem}[Poincar\'e-Bendixson]
\label{thm:poincare-bendixson}
Let $\dot{x} = f(x)$ be a $C^1$ system on an open subset of $\R^2$. Let
$\gamma^+(x_0)$ be a bounded positive semi-orbit contained in a compact set
$K$ with no equilibrium points. Then $\omega(x_0)$ is a periodic orbit
(and hence a limit cycle if $x_0 \notin \omega(x_0)$).
\end{theorem}

\begin{remark}
More generally, if $K$ contains finitely many equilibria, the $\omega$-limit
set $\omega(x_0)$ is either:
\begin{enumerate}
    \item an equilibrium point;
    \item a periodic orbit;
    \item a set consisting of equilibria connected by heteroclinic or
    homoclinic connections.
\end{enumerate}
\end{remark}

\begin{warning}[title=\textbf{Dimension 2 only}]
The Poincar\'e-Bendixson theorem is specific to dimension 2. In dimension
$n \geq 3$, bounded $\omega$-limit sets can be far more complicated (strange
attractors, for instance).
\end{warning}

\begin{corollary}[Existence of limit cycles]
Let $K \subset \R^2$ be a compact, positively invariant, annular set (without
holes) containing no equilibrium points. Then $K$ contains at least one
limit cycle.
\end{corollary}

\begin{intuition}
Strategy for proving the existence of a limit cycle:
\begin{enumerate}
    \item Construct an annular region $K$ that is positively invariant (the
    vector field points inward on the boundary).
    \item Verify that there are no equilibrium points in $K$.
    \item Conclude by Poincar\'e-Bendixson.
\end{enumerate}
\end{intuition}

\section{Poincar\'e Map and Multipliers}

\begin{definition}[Poincar\'e section]
A \textbf{Poincar\'e section} is a curve $\Sigma$ transverse to the vector
field, i.e., $f(x) \cdot n(x) \neq 0$ for all $x \in \Sigma$, where $n$
is a normal vector to $\Sigma$.
\end{definition}

\begin{definition}[Poincar\'e map]
The \textbf{Poincar\'e map} (or first return map) $P : \Sigma \to \Sigma$
maps a point $x \in \Sigma$ to the next intersection of its orbit with
$\Sigma$.
\end{definition}

\begin{proposition}
A periodic orbit $\Gamma$ corresponds to a fixed point of $P$. The stability
of $\Gamma$ is determined by $\abs{P'(x^*)}$:
\begin{itemize}
    \item $\abs{P'(x^*)} < 1$: stable cycle;
    \item $\abs{P'(x^*)} > 1$: unstable cycle;
    \item $\abs{P'(x^*)} = 1$: indeterminate case.
\end{itemize}
\end{proposition}

\begin{definition}[Floquet multiplier]
For a periodic orbit of period $T$ of the system $\dot{x} = f(x)$ in
dimension $n$, the \textbf{Floquet multipliers} are the eigenvalues of the
monodromy matrix
\[
    M = \Phi(T),
\]
where $\Phi(t)$ solves $\dot{\Phi} = Df(\gamma(t))\,\Phi$, $\Phi(0) = I$.
One multiplier always equals 1 (tangent direction). The orbit is stable if
all other multipliers have modulus $< 1$.
\end{definition}

\section{Index Theory}

\begin{definition}[Index of a closed curve]
Let $\gamma$ be a simple closed curve in $\R^2$ not passing through any
equilibrium. The \textbf{index of $\gamma$} with respect to the field $f$ is
the number of turns made by $f(x)$ as $x$ traverses $\gamma$ counterclockwise:
\[
    I_\gamma = \frac{1}{2\pi}\oint_\gamma d\theta, \qquad
    \theta = \arctan\frac{f_2(x)}{f_1(x)}.
\]
\end{definition}

\begin{theorem}[Properties of the index]
\begin{enumerate}
    \item If $\gamma$ encloses no equilibria, $I_\gamma = 0$.
    \item The index of a node or focus is $+1$.
    \item The index of a saddle is $-1$.
    \item The index of a periodic orbit is $+1$.
    \item The index is additive: if $\gamma$ encloses several equilibria,
    $I_\gamma = \sum I_{x_i^*}$.
\end{enumerate}
\end{theorem}

\begin{corollary}
Inside a limit cycle, the sum of indices of equilibria equals $+1$. In
particular, a limit cycle must enclose at least one equilibrium.
\end{corollary}

\section{Detailed Examples}

\begin{example}[Van der Pol oscillator --- existence of the cycle]
Consider the system in Li\'enard coordinates:
\[
    \dot{x} = y - F(x), \qquad \dot{y} = -x,
\]
with $F(x) = \mu(x^3/3 - x)$, $\mu > 0$. The unique equilibrium is the
origin (unstable focus for $\mu > 0$). We construct an annular region:
\begin{itemize}
    \item Inner boundary: small circle of radius $r$ (field points outward
    since the equilibrium is unstable);
    \item Outer boundary: energy level curve chosen large enough (the system
    is globally dissipative).
\end{itemize}
By Poincar\'e-Bendixson, a limit cycle exists in this region.
\end{example}

\begin{codeblock}[title=\textbf{Van der Pol limit cycle}]
\begin{minted}{python}
import numpy as np
from scipy.integrate import solve_ivp
import matplotlib.pyplot as plt

def van_der_pol(t, state, mu=1.0):
    x, y = state
    return [y, mu*(1 - x**2)*y - x]

fig, axes = plt.subplots(1, 3, figsize=(15, 5))
for ax, mu in zip(axes, [0.5, 1.0, 5.0]):
    for r in [0.1, 0.5, 1.0, 3.0, 5.0]:
        for theta in np.linspace(0, 2*np.pi, 4, endpoint=False):
            x0 = [r*np.cos(theta), r*np.sin(theta)]
            sol = solve_ivp(van_der_pol, [0, 50], x0,
                            args=(mu,), max_step=0.05)
            ax.plot(sol.y[0], sol.y[1], 'b-', lw=0.3, alpha=0.5)
    ax.set_title(f"$\\mu = {mu}$")
    ax.set_xlabel("$x$"); ax.set_ylabel("$\\dot{x}$")
    ax.set_xlim(-6, 6); ax.set_ylim(-8, 8)
    ax.set_aspect('auto'); ax.grid(True, alpha=0.3)
plt.tight_layout()
plt.savefig("van_der_pol_cycles.pdf")
plt.show()
\end{minted}
\end{codeblock}

\begin{example}[FitzHugh-Nagumo model]
The simplified FitzHugh-Nagumo neuron model is
\[
    \dot{v} = v - \frac{v^3}{3} - w + I, \qquad
    \dot{w} = \epsilon(v + a - bw),
\]
with $\epsilon \ll 1$. This system possesses a limit cycle for certain values
of the injected current $I$, modeling the action potential.
\end{example}

\section{Li\'enard's Theorem}

\begin{theorem}[Li\'enard]
Consider the Li\'enard equation $\ddot{x} + f(x)\dot{x} + g(x) = 0$ with
$F(x) = \int_0^x f(s)\, ds$. Under the conditions:
\begin{enumerate}
    \item $f$ and $g$ continuous, $g$ odd with $xg(x) > 0$ for $x \neq 0$;
    \item $F$ odd, $F(x) \to +\infty$ as $x \to +\infty$;
    \item $F$ is negative for $0 < x < a$, positive and increasing for
    $x > a$ (for some $a > 0$);
\end{enumerate}
then the system has a \textbf{unique stable limit cycle}.
\end{theorem}

\section{Exercises}

\begin{exercise}[Bendixson]
Show that the system $\dot{x} = x + e^{-y}$, $\dot{y} = y + e^x$ has no
periodic orbit in $\R^2$.
\end{exercise}

\begin{exercise}[Poincar\'e-Bendixson]
Prove the existence of a limit cycle for the system
$\dot{r} = r(1-r^2) + \epsilon r\cos\theta$,
$\dot{\theta} = 1$ for sufficiently small $\epsilon$.
\end{exercise}

\begin{exercise}[Index]
Compute the index of the origin for the system
$\dot{x} = x^2 - y^2$, $\dot{y} = 2xy$.
\end{exercise}

\begin{exercise}[Poincar\'e map]
For the system in polar coordinates
$\dot{r} = r(1-r)(2-r)$, $\dot{\theta} = 1$, find all limit cycles and
determine their stability.
\end{exercise}

\begin{exercise}[Van der Pol]
Show by the averaging method that for $\mu \ll 1$, the limit cycle of the
Van der Pol oscillator is close to the circle of radius 2.
\end{exercise}

\begin{exercise}[Uniqueness]
Let $\dot{x} = y$, $\dot{y} = -x + y(1 - x^2 - 3y^2)$. Show that there
exists a unique limit cycle using Li\'enard's theorem.
\end{exercise}
